\metatitle{Rowmotion and promotion}
\metadescription{Rowmotion and promotion on finite posets, with toggles,
illustrated examples, cyclic sieving, homomesy, and Tamari lattices.}
\metakeywords{rowmotion,promotion,posets,linear extensions,order ideals,
toggle group,cyclic sieving,homomesy,Tamari lattice}

\section[rowmotionPromotion]{Rowmotion and promotion}

Rowmotion and promotion are invertible operations on finite combinatorial
sets.  Promotion acts naturally on
\hyperref[linearExtension]{linear extensions}, while rowmotion acts on
\hyperref[orderIdeal]{order ideals}.  Both operations can be expressed as
products of local involutions.  This toggle description explains why their
orbits often admit rotation models, cyclic sieving, and homomesy.

Throughout, let $P$ be a finite poset with $n$ elements.  We write
$\mathcal L(P)$ for its linear extensions and $J(P)$ for its order ideals.

\section[posetPromotion]{Promotion on linear extensions}

Let $w=(p_1,p_2,\dotsc,p_n)\in\mathcal L(P)$.  For $1\leq i<n$, define an
adjacent involution $\tau_i$ by interchanging $p_i$ and $p_{i+1}$ when these
two elements are incomparable, and fixing $w$ otherwise.  The
\defin{promotion operator} is
\[
  \defin{\promotion}
  \coloneqq
  \tau_{n-1}\circ\tau_{n-2}\circ\dotsm\circ\tau_1,
\]
where $\tau_1$ is applied first.  This is Schützenberger promotion on linear
extensions \cite{Schutzenberger1972,Stanley2009}.

Equivalently, label each $p\in P$ by its position in $w$.  Remove the label
$1$, repeatedly slide down the smallest label among the elements covering
the empty position, place $n+1$ at the last position, and subtract $1$ from
every label.  The elements visited by the empty position form the
\defin{promotion chain}.

\begin{example*}[Promotion on a product of chains]
Let $P=[2]\times[3]$, ordered coordinatewise, and name its elements
\[
 a=(1,1),\quad b=(1,2),\quad c=(1,3),\quad
 d=(2,1),\quad e=(2,2),\quad f=(2,3).
\]
Its five linear extensions split into two promotion orbits:
\[
 abcdef\longmapsto abdecf\longmapsto adbcef\longmapsto abcdef,
 \qquad
 abdcef\longleftrightarrow adbecf.
\]
The vertex labels below give the positions of the elements in each word; the
three-cycle is in the top row and the two-cycle is in the bottom row.
\begin{figure}
\svgimg[width=0.66\textwidth]{svg-images/rowmotion-promotion-linear-extensions.svg}{
The two promotion orbits of linear extensions of the two-by-three product of
chains, shown as five labeled Hasse diagrams.}
\end{figure}
\end{example*}

For a Young-diagram poset, linear extensions are standard Young tableaux,
and this definition agrees with the
\hyperref[tableauPromotion]{tableau promotion operator}.  The convention for
promotion is not uniform in the literature; some authors use its inverse.

% Related Rust: rust/combinatoric-core/src/poset.rs
The methods \texttt{promotion\_linear\_extension},
\texttt{promotion\_orbit}, and \texttt{promotion\_orbits} in the shared Rust
library compute this action exactly.

\section[orderIdealToggles]{Order-ideal toggles}

For $p\in P$, the \defin{toggle at $p$}\label{orderIdealToggle} is the
involution $t_p:J(P)\to J(P)$ defined as follows.  If $p\notin I$, add $p$
when $I\cup\{p\}$ is an order ideal.  If $p\in I$, remove $p$ when
$I\setminus\{p\}$ is an order ideal.  In every other case, leave $I$ fixed.
The subgroup of permutations of $J(P)$ generated by the $t_p$ is the
\defin{toggle group} $T(P)$ \cite{StrikerWilliams2012}.

Two toggles commute unless their elements are joined by a cover relation.
Thus many global actions can be computed in several different orders without
changing the result.

\section[rowmotion]{Rowmotion}

The \defin{rowmotion operator} sends an order ideal to the ideal generated
by the minimal elements outside it:
\[
  \defin{\operatorname{Row}(I)}
  \coloneqq
  \bigl\{y\in P:y\leq_P x
  \text{ for some }x\in\min(P\setminus I)\bigr\}.
\]
If $(p_1,p_2,\dotsc,p_n)$ is any linear extension of $P$, rowmotion is also
obtained by toggling from maximal elements toward minimal elements.  More
precisely,
\[
  \operatorname{Row}
  =t_{p_1}\circ t_{p_2}\circ\dotsm\circ t_{p_n},
\]
so $t_{p_n}$ is applied first.  In particular, rowmotion is a bijection
\cite{StrikerWilliams2012}.

\begin{example*}[Two rowmotion orbits]
Return to the poset $a,b\lt c$.  Blue vertices in the figure are the elements
of an order ideal.  The five ideals split into an orbit of size three and an
orbit of size two:
\[
  \emptyset\longmapsto\{a,b\}\longmapsto\{a,b,c\}
  \longmapsto\emptyset,
  \qquad
  \{a\}\longleftrightarrow\{b\}.
\]
\begin{figure}
\svgimg[width=0.78\textwidth]{svg-images/rowmotion-promotion-order-ideals.svg}{
The two rowmotion orbits of order ideals of the V-shaped poset.}
\end{figure}
For instance, the minimal elements outside $\{a\}$ consist only of $b$, so
the generated ideal is $\{b\}$.
\end{example*}

The shared Rust library implements the same definition through
\texttt{rowmotion\_order\_ideal}, \texttt{rowmotion\_orbit}, and
\texttt{rowmotion\_orbits}.

\section[rowmotionPromotionConjugacy]{Promotion and rowmotion through toggles}

There is also a toggle version of promotion on order ideals.  Suppose that a
poset is drawn with rows and columns so that every cover joins diagonally
adjacent positions.  Toggling one column at a time defines promotion, while
toggling one row at a time defines rowmotion.

\begin{theorem}[Striker--Williams, \cite{StrikerWilliams2012}]
On every rowed-and-columned poset, toggle promotion and rowmotion are
conjugate in the toggle group.  Consequently, they have the same orbit
structure.
\end{theorem}

In the Ferrers and root-poset examples, boundary paths convert the
column-by-column action into familiar promotion or rotation.  This gives, for
example, an equivariant bijection between rowmotion on $J([a]\times[b])$ and
rotation of binary words with $a$ ones and $b$ zeros.  Hence rowmotion has
order $a+b$ on this product of chains, and
\[
  \left(J([a]\times[b]),\langle\operatorname{Row}\rangle,
  \qbinom{a+b}{a}_q\right)
\]
exhibits the \hyperref[cyclicSieving]{cyclic sieving phenomenon}.

Other rotation models connect rowmotion on positive root posets with
Kreweras complementation on noncrossing partitions.  See the cyclic-sieving
sections on \hyperref[cspListPosetsRoot]{root posets} and
\hyperref[cspListPosetsMinuscule]{minuscule posets}.  Orbit averages of
statistics under rowmotion are among the main examples of
\hyperref[homomesy]{homomesy}.

\section[rowmotionTamari]{Rowmotion on Tamari lattices}

The classical definition above treats the distributive lattice $J(P)$.
Rowmotion also extends to semidistributive lattices through their canonical
edge labels.  Alt $\nu$-Tamari lattices form a family of such lattices on
$\nu$-Dyck paths; for a fixed path $\nu$, the family interpolates between a
Tamari-type lattice and the distributive nesting lattice.

\begin{theorem}[Adenbaum--Barnard--Ceballos--Chenevière--Defant--Hopkins--Müller--Rubey--Striker,
\cite[Thm.~1.1]{AdenbaumBarnardCeballosCheneviereDefantHopkinsMullerRubeyStriker2026x}]
Fix a lattice path $\nu$.  Any two alt $\nu$-Tamari lattices admit a bijection
that intertwines rowmotion and preserves the orbit sums of the down-degree
statistic.
\end{theorem}

Thus rowmotion has the same orbit structure throughout the family, even
though the underlying lattices need not be isomorphic.  The theorem transfers
two particularly clean statements from distributive lattices to all rational
alt Tamari lattices
\cite[Thm.~5.1--5.2]{AdenbaumBarnardCeballosCheneviereDefantHopkinsMullerRubeyStriker2026x}:

\begin{itemize}
\item If $a\lt b$ are coprime, rowmotion on every rational $(a,b)$ alt Tamari
lattice has order dividing $a+b-1$, and its orbit structure is encoded by the
rational $q$-Catalan polynomial
\[
  \frac{1}{[a+b]_q}\qbinom{a+b}{a}_q.
\]
\item On every alt $m$-Tamari lattice of size parameter $n$, the down-degree
statistic is homomesic with average $m(n-1)/(m+1)$.
\end{itemize}

It remains open whether the down-degree statistic on every rational
$(a,b)$ alt Tamari lattice has orbit average
$(a-1)(b-1)/(a+b-1)$.  The authors also conjecture that rowmotion invariance
extends from alt $\nu$-Tamari lattices to cross Tamari lattices associated
with moon polyominoes.
